\scanpage{52}
\unsure{vu que l'opération} \illegible{1} \unsure{de passage au sous-complexe des}
\unsure{cochaînes non dégénérées commute à}
$\operatorname{Hom}(X_*,-)$.
Donc
$\unsure{C^{\prime *}(X_*)=\operatorname{Hom}(X_*,C_{*M}^{\prime *})}$
\unsure{satisfait encore aux conditions} a) et b).
\unsure{D'autre part, les} $\unsure{C_{*M}^{\prime n}}$ \unsure{sont flasques},
\illegible{2} \unsure{explicite} \unsure{aisément}
\[
\begin{aligned}
\unsure{(C_{*M}^{\prime n})_m}
&\simeq M^{\unsure{\operatorname{Hom}_{\mathrm{mono}}(\Delta_n,\Delta_m)}}
\simeq k^{\unsure{\operatorname{Hom}_{\mathrm{mono}}(\Delta_n,\Delta_m)}}\otimes_k M\\
&\simeq \bigl(\textstyle\bigwedge^{n+1}k^{\Delta_m}\bigr)\otimes_k M.
\end{aligned}
\]
\unsure{et cette formule} \illegible{4}, d'où
\[
\unsure{C_{*M}^{\prime n}\simeq\bigwedge^{n+1}Q_{*k}\otimes_k M}
\]
où $\unsure{Q_{*k}=(\Delta_*\longmapsto k^{\Delta_*})}$.
[\unsure{les} \illegible{2} \unsure{des} \illegible{1}

\unsure{Soit} \illegible{2}
\[
\unsure{T:k_*\longrightarrow Q_{*k}=C_{*k}^{\prime0}}
\]
\illegible{2} \unsure{est la coaugmentation}
($\unsure{Q_{-1,k}=k}$).

\unsure{Alors} \illegible{2} \unsure{s'identifie},
\unsure{donc}
$\unsure{C_{*M}^{\prime n}\simeq\bigwedge^{n+1}Q_{*k}\otimes_k M}$
\unsure{étant muni} \unsure{par} $\unsure{T\wedge}$.
\unsure{C'est} \illegible{3} \unsure{pour calculer les} \illegible{1}].
\unsure{Or} \illegible{2} \unsure{lorsque} \illegible{2}
\unsure{particulier}, \illegible{3},
$\unsure{\bigwedge^i Q_{*k}\otimes_k M}$.
\illegible{3} $\unsure{\bigwedge^0Q_{*k}\otimes_k M}$,
\illegible{3} \unsure{de} \illegible{2}.

\unsure{On retrouve en même temps la} \illegible{2}
$\unsure{C_{*k}^{\prime *}}$:
\unsure{c'est un complexe de} \illegible{2},
\illegible{3}
\unsure{dont} \illegible{1} (\unsure{savoir pour les}
\illegible{2} \unsure{les})
\unsure{une structure} \illegible{2}
$\unsure{\partial:C_{*M}^n\longrightarrow C_{*M}^{n+1}}$
\illegible{2} (\unsure{sauf erreur}).
\mspage{51}
